% !PRINTED 306
\weberpage{306}

Wir machen beim Beweis dieses Satzes die folgenden Voraussetzungen, von
denen wir uns zum Theil später wieder befreien werden.

1. $f(x)$ hat keine doppelten oder mehrfachen Wurzeln.

2. Von den Functionen $f_\nu(x)$ verschwinden nicht zwei benachbarte für
denselben Werth von $x$.

3. Auch von den Functionen $F_\nu(x)$ verschwinden nicht zwei benachbarte
für denselben Werth von $x$.

4. Eine Folge von 2. ist [nach Formel (3)], dass nicht $f_\nu(x)$ und
$F_\nu(x)$ für denselben Werth von $x$ verschwinden.

5. Für $x = \alpha$ und $x = \beta$ selbst soll keine der Functionen
$f_\nu$, $F_\nu$ verschwinden.

Der Beweis der Sätze II und III ist nun folgender:

Lassen wir $x$, stetig wachsend, von $\alpha$ bis $\beta$ gehen, so kann
eine Aenderung in der Zeichenfolge nur dann eintreten, wenn eine der
Functionen $f_\nu$ oder $F_\nu$ durch Null hindurchgeht, und es kommt nur
darauf an, den Effect zu untersuchen, den ein solcher Durchgang durch
Null in den verschiedenen Fällen hervorruft. Nehmen wir zunächst an, es
gehe eine der mittleren Functionen $f$, etwa $f_\nu$, durch Null, wobei
also $0 < \nu < n$. Je nachdem $f_{\nu+1}$ positiv oder negativ ist,
wird $f_\nu$ wachsen oder abnehmen. Es wird also im ersten Falle
$f_\nu(x - 0)$ negativ, $f_\nu(x + 0)$ positiv sein, und im zweiten
Falle umgekehrt; ausserdem folgt aus (3), dass, wenn $f_\nu = 0$ ist,
$F_{\nu-1}$ und $F_{\nu+1}$ positiv sind, und dass $F_\nu$ das
entgegengesetzte Vorzeichen von $f_{\nu-1} f_{\nu+1}$ hat. Wir stellen
die hiernach bleibenden Möglichkeiten tabellarisch zusammen, lassen aber
in der Aufzählung solche Fälle weg, die durch gleichzeitige
Vorzeichenänderung der drei Functionen $f_{\nu-1}$, $f_\nu$,
$f_{\nu+1}$ aus den aufgeführten entstehen, da diese offenbar keine
verschiedenen Resultate ergeben. Es bleiben uns hiernach nur zwei
Fälle:

% [restored 2026-08-07, 6x zoom strips 5/6 and 4/4] The lower third of this page was dropped:
% two side-by-side sign tables, the connecting line "In diesen Fällen ist also", and a second
% pair of tables. The fragment ended at "nur zwei Fälle:" -- announcing tables that were not
% there -- and carried no %CONT>> marker, so the material was not deferred either. i0332 opens
% with "Es wird also weder ein $VP$ noch ein $PP$ verloren oder gewonnen.", which is printed 307's
% own first sentence, so nothing was moved: it was lost.
%
% Not caught by any mechanical check. Coverage did not move because a sign table is almost no
% letters, and the fragment was internally well-formed. Raised as an omission by this page's
% translation pass.
%
% The printed tables carry outer vertical rules as well as the rules separating the row labels
% and the F row; all are reproduced. The heading of each half repeats "\nu-1, \nu, \nu+1".
\[
\begin{array}{|c|ccc|}
 & \nu-1, & \nu, & \nu+1 \\ \hline
f(x-0) & + & - & + \\
f(x+0) & + & + & + \\ \hline
F & + & - & +
\end{array}
\qquad
\begin{array}{|ccc|}
\nu-1, & \nu, & \nu+1 \\ \hline
- & - & + \\
- & + & + \\ \hline
+ & + & +
\end{array}.
\]
In diesen Fällen ist also
\[
\begin{array}{|c|cc|}
x-0 & VV, & VV \\
x+0 & PV, & PV
\end{array}
\qquad
\begin{array}{|cc|}
PP, & VP \\
VP, & PP
\end{array}.
\]
